\chapter{Morphisms and Quotients}

\begin{center}\itshape
Computing modulo $n$ forgets everything about an integer except its remainder, and the arithmetic still works.\\
The sign of a permutation forgets everything except its parity, and products still behave.\\
Both are the same construction: a quotient group.
\end{center}

The integers $4$, $7$ and $10$ are different. But for a clock that records only the remainder on division by $3$, all three read $1$. Addition still works after we forget the multiples of $3$:
\[
[4]+[7]=[11]=[2]\quad\text{in }\Z/3\Z.
\]
We can calculate with the shorter descriptions because forgetting the information is compatible with addition.

A group homomorphism works the same way. It sends elements of one group to another while preserving products. Its \emph{kernel} records exactly what has been forgotten. For reduction modulo $3$, the kernel is $3\Z$. Every integer in a coset of $3\Z$ gives the same remainder, and different cosets give different remainders.

Not every subgroup can be treated as a set of elements to be collapsed. If we try to multiply cosets using arbitrary representatives, the answer may change when we choose a different representative. Normality is the condition that stops this from happening. It gives a group $G/N$, the quotient of $G$ by $N$.

This matters beyond clocks. The sign of a permutation forgets every detail except whether it is even or odd. The determinant of an invertible matrix forgets its shape but remembers the factor by which it scales oriented volume. Both maps have kernels and both lead to useful quotient groups. And when we want to build a complicated group from simpler ones, products and semidirect products let us reverse the process.

\section{Homomorphisms}

\begin{definition}
A \emph{homomorphism} $f:G\to G'$ is a map with $f(ab)=f(a)f(b)$ for all $a,b\in G$. Its \emph{kernel} is $\ker f=\{g\in G:f(g)=e'\}$ and its \emph{image} is $\operatorname{im}f=f(G)$.
\end{definition}

As for isomorphisms (Chapter~1), $f(e)=e'$, $f(g^{-1})=f(g)^{-1}$, and hence $f(g^k)=f(g)^k$ for all $k\in\Z$.

\begin{example}\label{ex:homs}
\begin{enumerate}[label=(\arabic*)]
\item Reduction $\Z\to\Z/n\Z$, $k\mapsto k\bmod n$, with kernel $n\Z$.
\item The determinant $\det:GL_n(K)\to K^*$, with kernel $SL_n(K)$ (Volume~1, Chapter~8: $\det(AB)=\det A\det B$).
\item The sign $\varepsilon:\mathfrak S_n\to\{\pm1\}$ (Volume~1, Chapter~8, the sign lemma), with kernel the \emph{alternating group} $A_n$ of even permutations.
\item The exponential $(\R,+)\to(\R_{>0},\times)$, an isomorphism; and $x\mapsto e^{2i\pi x}$ from $(\R,+)$ onto the circle $\mathbb U=\{z\in\C:|z|=1\}$, with kernel $\Z$.
\item For any $g\in G$, the \emph{conjugation} $c_g:x\mapsto gxg^{-1}$ is an isomorphism $G\to G$ (an \emph{automorphism}): $gxyg^{-1}=(gxg^{-1})(gyg^{-1})$, and its inverse is $c_{g^{-1}}$.
\end{enumerate}
\end{example}

\begin{proposition}\label{prop:kerim}
Let $f:G\to G'$ be a homomorphism.
\begin{enumerate}[label=(\alph*)]
\item $\ker f$ is a subgroup of $G$ and $\operatorname{im}f$ a subgroup of $G'$.
\item $f(a)=f(b)$ if and only if $a\in b\ker f$. In particular $f$ is injective if and only if $\ker f=\{e\}$.
\item For $g\in G$ and $k\in\ker f$, $gkg^{-1}\in\ker f$.
\end{enumerate}
\end{proposition}
\begin{proof}
(a) If $f(a)=f(b)=e'$ then $f(ab^{-1})=f(a)f(b)^{-1}=e'$; and $e\in\ker f$. For the image, $f(a)f(b)^{-1}=f(ab^{-1})$. Apply Proposition~\ref{prop:subgroup}(a) both times.
(b) $f(a)=f(b)$ iff $f(b^{-1}a)=f(b)^{-1}f(a)=e'$ iff $b^{-1}a\in\ker f$ iff $a\in b\ker f$. If $\ker f=\{e\}$ this says $a=b$.
(c) $f(gkg^{-1})=f(g)\,e'\,f(g)^{-1}=e'$.
\end{proof}

Part (b) says that the fibres of $f$, the sets $f^{-1}(y)$, are exactly the cosets of $\ker f$. Here is the familiar case of reduction modulo $3$. The numbers $-3$, $0$ and $3$ have different values in $\Z$, but the same image in $\Z/3\Z$. The map has forgotten multiples of $3$ and nothing else.

\begin{figure}[ht]
\centering
\fig{ch02-kernel-fibres}
\caption{Reduction modulo $3$ groups integers into three fibres. Each fibre is a coset of the kernel $3\Z$.}
\label{fig:kernel-fibres}
\end{figure}

The picture has only nine source integers, but each fibre continues forever in both directions. In the next theorem, a fibre becomes a single element of the quotient group.

A kernel is a subgroup, but it has a special property: conjugating one of its elements by any group element keeps us inside the kernel. This is exactly the condition we need when we try to multiply entire cosets instead of individual elements.

\section{Normal subgroups and quotients}

\begin{think}
Look again at $\mathfrak S_3$. The subgroup $A_3=\{e,(1\,2\,3),(1\,3\,2)\}$ is the kernel of the sign. Could the subgroup $H=\{e,(1\,2)\}$ be the kernel of some homomorphism from $\mathfrak S_3$? Compare $(1\,3)H$ with $H(1\,3)$, and use Proposition~\ref{prop:kerim}(c).
\end{think}

\begin{definition}
A subgroup $N\le G$ is \emph{normal}, written $N\trianglelefteq G$, if $gNg^{-1}\subseteq N$ for every $g\in G$.
\end{definition}

\begin{proposition}\label{prop:normal}
For a subgroup $N\le G$, the following are equivalent:
\begin{enumerate}[label=(\roman*)]
\item $N$ is normal;
\item $gNg^{-1}=N$ for every $g\in G$;
\item $gN=Ng$ for every $g\in G$: left and right cosets coincide.
\end{enumerate}
Kernels of homomorphisms are normal; every subgroup of an abelian group is normal; and a subgroup of index $2$ is normal.
\end{proposition}
\begin{proof}
(i)$\Rightarrow$(ii): apply (i) to $g^{-1}$: $g^{-1}Ng\subseteq N$, so $N\subseteq gNg^{-1}$. (ii)$\Rightarrow$(iii): multiply $gNg^{-1}=N$ on the right by $g$. (iii)$\Rightarrow$(i): $gN g^{-1}=Ngg^{-1}=N$.
Kernels are normal by Proposition~\ref{prop:kerim}(c), and in an abelian group $gng^{-1}=n$. If $[G:N]=2$, the left cosets are $N$ and $G\setminus N$, and so are the right cosets (the right cosets also partition $G$ into sets of size $|N|$, by the same proof as Lagrange's theorem). So for $g\notin N$, $gN=G\setminus N=Ng$, and for $g\in N$, $gN=N=Ng$.
\end{proof}

In $\mathfrak S_3$, $H=\{e,(1\,2)\}$ is not normal: $(1\,3)(1\,2)(1\,3)=(2\,3)\notin H$. So $H$ is not a kernel. The rotations of $D_n$, of index $2$, are normal, and so is $A_n$ in $\mathfrak S_n$. Figure~\ref{fig:lattices} shows which subgroups of $D_4$ are normal.

\begin{theorem}[quotient group]\label{thm:quotient}
Let $N\trianglelefteq G$, and let $G/N$ be the set of cosets $gN$. The rule
\[
(aN)(bN)=abN
\]
is well defined and makes $G/N$ a group, with identity $N$ and $(aN)^{-1}=a^{-1}N$. The projection $\pi:G\to G/N$, $g\mapsto gN$, is a surjective homomorphism with kernel $N$. If $G$ is finite, $|G/N|=|G|/|N|$.

Conversely, if $N\le G$ and the rule above is well defined, then $N$ is normal.
\end{theorem}
\begin{proof}
Well defined: suppose $aN=a'N$ and $bN=b'N$, so $a'=an$ and $b'=bm$ with $n,m\in N$. Then $a'b'=anbm=ab\,(b^{-1}nb)\,m$, and $b^{-1}nb\in N$ by normality, so $a'b'\in abN$ and $a'b'N=abN$. The group axioms for $G/N$ follow from those of $G$: $(aN\,bN)\,cN=(ab)cN=a(bc)N=aN\,(bN\,cN)$, $N\,aN=aN$, and $aN\,a^{-1}N=N$. By definition $\pi(ab)=\pi(a)\pi(b)$, $\pi$ is onto, and $\pi(g)=N$ iff $g\in N$. The order is $[G:N]=|G|/|N|$ by Lagrange.
Conversely, if the rule is well defined then for $g\in G$, $n\in N$: $gnN=gN$ (as $nN=N$), and multiplying by $g^{-1}N$ gives $gng^{-1}N=gg^{-1}N=N$, i.e.\ $gng^{-1}\in N$.
\end{proof}

\paragraph{Why this works.}
There is one place where the quotient needs a genuine proof: we may write the same coset as $aN$ or as $anN$ with $n\in N$. If we choose another representative of the second coset as well, the extra element $n$ has to pass through $b$. Normality guarantees that $b^{-1}nb$ is still in $N$, so it disappears when we pass to cosets. Without that step, $(aN)(bN)$ could depend on which representatives we happened to pick. Once independence is settled, the group laws are inherited from $G$.

So the normal subgroups are exactly the subgroups by which one can divide, and they are exactly the kernels: $N$ is the kernel of $\pi$.

\begin{theorem}[first isomorphism theorem]\label{thm:iso1}
Let $f:G\to G'$ be a homomorphism. Then $\bar f:G/\ker f\to\operatorname{im}f$, $g\ker f\mapsto f(g)$, is a well-defined isomorphism. In particular, for finite $G$, $|G|=|\ker f|\cdot|\operatorname{im}f|$.
\end{theorem}
\begin{proof}
Write $K=\ker f$. By Proposition~\ref{prop:kerim}, $K$ is normal, so $G/K$ is a group. We would like to send the coset $gK$ to the element $f(g)$. Before doing anything else, we must check that the answer does not depend on the representative $g$.

Suppose $gK=hK$. Then $h^{-1}g\in K$, so
\[
f(h)^{-1}f(g)=f(h^{-1}g)=e.
\]
Thus $f(g)=f(h)$: the rule $\bar f(gK)=f(g)$ is well defined. Conversely, if $f(g)=f(h)$, the same equation gives $h^{-1}g\in K$, hence $gK=hK$. This converse makes $\bar f$ injective.

For multiplication, use the definition of the quotient product:
\[
\bar f((gK)(hK))=\bar f(ghK)=f(gh)=f(g)f(h)
=\bar f(gK)\bar f(hK).
\]
So $\bar f$ is a homomorphism. Its image is exactly $\operatorname{im}f$ by definition; viewed as a map to that image, it is surjective. Hence it is an isomorphism.

When $G$ is finite, Lagrange gives $|G/K|=|G|/|K|$. Since $G/K\cong\operatorname{im}f$, we obtain $|G|=|\ker f|\,|\operatorname{im}f|$.
\end{proof}

\paragraph{Why this works.}
A homomorphism makes two elements indistinguishable precisely when their quotient lies in the kernel: $f(a)=f(b)$ iff $f(b^{-1}a)=e$. The quotient $G/\ker f$ records one representative of each indistinguishable class, while $\operatorname{im}f$ records the distinct values. The theorem says these are the same group with two names. For the reduction map $\Z\to\Z/4\Z$, the kernel is $4\Z$; the theorem identifies the four cosets with the four residues, not with four arbitrary integers.

\begin{example}\label{ex:iso1}
\begin{enumerate}[label=(\arabic*)]
\item $\Z/n\Z$ is the quotient of $\Z$ by $n\Z$, as in Volume~2, and for $n\ge2$, $\mathfrak S_n/A_n\cong\{\pm1\}$: so $|A_n|=\frac{n!}2$.
\item $GL_n(K)/SL_n(K)\cong K^*$, since $\det$ is onto (use $\mathrm{diag}(\lambda,1,\dots,1)$).
\item $\R/\Z\cong\mathbb U$, by $x\mapsto e^{2i\pi x}$: the circle is the line rolled up. Its elements of finite order form $\Q/\Z$, the group of all roots of unity.
\item If $G=\langle g\rangle$ has order $n$, then $k\mapsto g^k$ is onto with kernel $n\Z$, which recovers $G\cong\Z/n\Z$ (Theorem~\ref{thm:cyclic}).
\end{enumerate}
\end{example}

\begin{theorem}[second and third isomorphism theorems, correspondence]\label{thm:iso23}
Let $N\trianglelefteq G$ and $\pi:G\to G/N$.
\begin{enumerate}[label=(\alph*)]
\item (Correspondence) $H\mapsto H/N=\pi(H)$ is a bijection from the subgroups of $G$ containing $N$ to the subgroups of $G/N$, which preserves inclusions and normality.
\item (Third theorem) If $N\le M\trianglelefteq G$, then $(G/N)/(M/N)\cong G/M$.
\item (Second theorem) If $H\le G$, then $HN=\{hn\}$ is a subgroup of $G$, $H\cap N\trianglelefteq H$, and $H/(H\cap N)\cong HN/N$.
\end{enumerate}
\end{theorem}
\begin{proof}
(a) The image $\pi(H)$ of a subgroup is a subgroup, and the preimage $\pi^{-1}(\bar H)$ of a subgroup $\bar H\le G/N$ is a subgroup containing $N=\pi^{-1}(\{N\})$ (preimages of subgroups under homomorphisms are subgroups, as in Proposition~\ref{prop:kerim}(a)). These two maps are inverse to each other: $\pi(\pi^{-1}(\bar H))=\bar H$ since $\pi$ is onto, and $\pi^{-1}(\pi(H))=HN=H$ when $N\subseteq H$. Both preserve inclusions. If $H\trianglelefteq G$ then $\pi(g)\pi(H)\pi(g)^{-1}=\pi(gHg^{-1})=\pi(H)$; if $\pi(H)\trianglelefteq G/N$ then $\pi(ghg^{-1})\in\pi(H)$ gives $ghg^{-1}\in HN=H$.
(b) The map $G/N\to G/M$, $gN\mapsto gM$, is well defined (as $N\subseteq M$), a homomorphism and onto; its kernel is $\{gN:g\in M\}=M/N$. Apply Theorem~\ref{thm:iso1}.
(c) For $h,h'\in H$ and $n,n'\in N$, $(hn)(h'n')^{-1}=hh'^{-1}\,(h'n n'^{-1}h'^{-1})\in HN$, so $HN$ is a subgroup; it contains $N$, which is normal in it. The map $H\to HN/N$, $h\mapsto hN$, is a homomorphism, onto since $hnN=hN$, with kernel $H\cap N$. Apply Theorem~\ref{thm:iso1}.
\end{proof}

\begin{figure}[ht]
\centering
\fig{ch02-lattices}
\caption{The subgroups of $D_4$ (left) and of $\Z/12\Z$ (right), each joined to the subgroups just above it that contain it. Normal subgroups are in charcoal. In $\Z/12\Z$ every subgroup is normal, there is one subgroup of each order $d\mid12$ (Theorem~\ref{thm:cyclic}), and the picture is the lattice of divisors of $12$. In $D_4$, $6$ of the $10$ subgroups are normal; the four reflection subgroups $\langle s\rangle,\dots$ are not, and $\langle s\rangle$ is normal in the Klein group $\langle s,sr^2\rangle$, which is normal in $D_4$, although $\langle s\rangle$ is not normal in $D_4$.}
\label{fig:lattices}
\end{figure}

Quotients take a group apart by identifying elements. We can also assemble groups, first by taking pairs and then by letting one group act on another. The second construction is what gives the dihedral group its mixture of rotations and reflections.

\section{Products}

\begin{definition}
The \emph{direct product} $G\times H$ is the set of pairs with $(g,h)(g',h')=(gg',hh')$.
\end{definition}

It is a group of order $|G||H|$, abelian if $G$ and $H$ are. The order of $(g,h)$ is $\operatorname{lcm}(o(g),o(h))$, since $(g,h)^k=(g^k,h^k)$ is trivial iff both $o(g)$ and $o(h)$ divide $k$.

\begin{theorem}[Chinese remainder theorem, group version]\label{thm:crt}
$\Z/mn\Z\cong\Z/m\Z\times\Z/n\Z$ if and only if $\gcd(m,n)=1$.
\end{theorem}
\begin{proof}
If $\gcd(m,n)=1$, the map $k\bmod mn\mapsto(k\bmod m,k\bmod n)$ is a well-defined homomorphism. Its kernel is the set of classes of $k$ with $m\mid k$ and $n\mid k$, i.e.\ $mn\mid k$: it is trivial, so the map is injective, hence bijective since both sides have $mn$ elements. If $d=\gcd(m,n)>1$, every element of $\Z/m\Z\times\Z/n\Z$ has order dividing $\operatorname{lcm}(m,n)=\frac{mn}d<mn$, so this group is not cyclic, while $\Z/mn\Z$ is.
\end{proof}

How can we recognise a product inside a group?

\begin{proposition}[internal direct product]\label{prop:internal}
Let $H,K$ be normal subgroups of $G$ with $H\cap K=\{e\}$ and $HK=G$. Then $hk=kh$ for $h\in H$, $k\in K$, and $(h,k)\mapsto hk$ is an isomorphism $H\times K\to G$.
\end{proposition}
\begin{proof}
The \emph{commutator} $hkh^{-1}k^{-1}$ lies in $K$ (as $hkh^{-1}\in K$) and in $H$ (as $kh^{-1}k^{-1}\in H$), so it is $e$: $hk=kh$. Then $(h,k)\mapsto hk$ is a homomorphism, since $hkh'k'=hh'kk'$. Its kernel is $\{(h,k):h=k^{-1}\}$, and $h=k^{-1}\in H\cap K$ forces $h=k=e$. It is onto because $HK=G$.
\end{proof}

When only one of the two subgroups is normal, the group is a twisted product. In $D_n$, the rotations $N=\langle r\rangle$ form a normal subgroup, $H=\langle s\rangle$ does not, $N\cap H=\{e\}$ and $NH=D_n$; the elements still factor uniquely as $r^ks^\varepsilon$, but $s$ acts on $N$ by inverting it.

\begin{definition}
Let $N,H$ be groups and $\varphi:H\to\operatorname{Aut}(N)$ a homomorphism, written $h\mapsto\varphi_h$. The \emph{semidirect product} $N\rtimes_\varphi H$ is the set $N\times H$ with the product
\[
(n,h)(n',h')=(n\,\varphi_h(n'),\ hh').
\]
\end{definition}

\begin{proposition}\label{prop:semidirect}
\begin{enumerate}[label=(\alph*)]
\item $N\rtimes_\varphi H$ is a group of order $|N||H|$, in which $N\times\{e\}$ is a normal subgroup isomorphic to $N$, $\{e\}\times H$ a subgroup isomorphic to $H$, and $(e,h)(n,e)(e,h)^{-1}=(\varphi_h(n),e)$.
\item (Recognition) If $N\trianglelefteq G$, $H\le G$, $N\cap H=\{e\}$ and $NH=G$, then $G\cong N\rtimes_\varphi H$ with $\varphi_h(n)=hnh^{-1}$.
\end{enumerate}
\end{proposition}
\begin{proof}
(a) Associativity: the two ways of bracketing give
\[\begin{aligned}{}&\big((n,h)(n',h')\big)(n'',h'')\eqbr\quad=\big(n\varphi_h(n')\varphi_{hh'}(n''),\ hh'h''\big),\\{}&(n,h)\big((n',h')(n'',h'')\big)\eqbr\quad=\big(n\varphi_h(n'\varphi_{h'}(n'')),\ hh'h''\big),\end{aligned}\]
and these agree because $\varphi_h$ is a homomorphism and $\varphi_h\circ\varphi_{h'}=\varphi_{hh'}$. The identity is $(e,e)$, and the inverse of $(n,h)$ is $(\varphi_{h^{-1}}(n^{-1}),h^{-1})$, as one checks. The map $(n,h)\mapsto h$ is a homomorphism onto $H$ with kernel $N\times\{e\}$, which is therefore normal. The last formula is a direct computation: $(e,h)(n,e)=(\varphi_h(n),h)$ and $(\varphi_h(n),h)(e,h^{-1})=(\varphi_h(n),e)$.
(b) Conjugation by $h$ preserves the normal subgroup $N$, so $\varphi$ is a homomorphism $H\to\operatorname{Aut}(N)$. The map $(n,h)\mapsto nh$ satisfies $nh\,n'h'=n(hn'h^{-1})hh'$, so it is a homomorphism from $N\rtimes_\varphi H$. It is onto as $NH=G$, and injective because $nh=e$ gives $n=h^{-1}\in N\cap H=\{e\}$.
\end{proof}

So $D_n\cong\Z/n\Z\rtimes\Z/2\Z$, where the generator of $\Z/2\Z$ acts by $k\mapsto-k$. With the trivial action, $\varphi_h=\mathrm{id}$, we get back the direct product.

\begin{beauty}
The first isomorphism theorem is a single sentence, yet it organises everything. To understand a group, find a homomorphism out of it: its kernel is a smaller group, its image another one, and the group is built from these two pieces. Repeating this until the pieces cannot be broken further leads to the simple groups, the ``prime numbers'' of group theory. We meet the first non-abelian one, $A_5$, in Chapter~4, and in Chapter~10 the fact that $A_5$ cannot be broken into abelian pieces becomes the reason why the quintic has no formula.
\end{beauty}

\section{Exercises}

\begin{exercise}[$\star$]\label{ex:autzn}
Show that every homomorphism $\Z/n\Z\to\Z/n\Z$ is $x\mapsto ax$ for a unique $a\in\Z/n\Z$, and deduce that $\operatorname{Aut}(\Z/n\Z)\cong(\Z/n\Z)^*$.
\end{exercise}

\begin{exercise}[$\star$]\label{ex:centre}
Show that the centre $Z(G)$ is a normal subgroup. Show that if $G/Z(G)$ is cyclic, then $G$ is abelian.
\end{exercise}

\begin{exercise}[$\star$]\label{ex:phimult}
Deduce from Theorem~\ref{thm:crt} that $(\Z/mn\Z)^*\cong(\Z/m\Z)^*\times(\Z/n\Z)^*$ for coprime $m,n$, and recover that $\varphi$ is multiplicative.
\end{exercise}

\begin{exercise}[$\star\star$]\label{ex:commutator}
The \emph{derived subgroup} $G'$ of $G$ is the subgroup generated by the commutators $[a,b]=aba^{-1}b^{-1}$. Show that $G'$ is normal, that $G/G'$ is abelian, and that for $N\trianglelefteq G$, $G/N$ is abelian if and only if $G'\subseteq N$. Compute $\mathfrak S_3'$.
\end{exercise}

\begin{exercise}[$\star\star$]\label{ex:quaternion}
\emph{The quaternion group.} Let $Q_8=\{\pm1,\pm i,\pm j,\pm k\}$ with $i^2=j^2=k^2=ijk=-1$ (realised by the matrices $\begin{pmatrix}i&0\\0&-i\end{pmatrix}$, $\begin{pmatrix}0&1\\-1&0\end{pmatrix}$ and their product in $GL_2(\C)$). Show that every subgroup of $Q_8$ is normal although $Q_8$ is not abelian, and that $Q_8$ is not a semidirect product of two proper subgroups.
\end{exercise}

\begin{exercise}[$\star\star$]\label{ex:affine}
\emph{The affine group.} Let $p$ be prime and $\mathrm{Aff}_p$ the group of maps $x\mapsto ax+b$ of $\Z/p\Z$, with $a\ne0$. Show that it has order $p(p-1)$, that the translations form a normal subgroup $T$, and that $\mathrm{Aff}_p\cong\Z/p\Z\rtimes(\Z/p\Z)^*$. Show that $\mathrm{Aff}_3\cong\mathfrak S_3$.
\end{exercise}

\begin{exercise}[$\star\star$]\label{ex:normalnormal}
Show, using Figure~\ref{fig:lattices}, that ``being a normal subgroup of'' is not transitive. Show that a subgroup $K$ that is \emph{characteristic} in $N$ (fixed by every automorphism of $N$) and $N\trianglelefteq G$ imply $K\trianglelefteq G$. Deduce that every subgroup of a normal cyclic subgroup is normal.
\end{exercise}

\begin{exercise}[$\star\star\star$]\label{ex:order8}
\emph{The groups of order $8$.} Show that every group of order $8$ is isomorphic to exactly one of $\Z/8\Z$, $\Z/4\Z\times\Z/2\Z$, $(\Z/2\Z)^3$, $D_4$, $Q_8$.
\end{exercise}

\section{Hints}

\begin{hint}{ex:autzn}
A homomorphism from a cyclic group is determined by the image $a$ of a generator. When is $x\mapsto ax$ bijective?
\end{hint}
\begin{hint}{ex:centre}
If $G/Z(G)$ is generated by $gZ(G)$, every element is $g^kz$ with $z$ central.
\end{hint}
\begin{hint}{ex:phimult}
An isomorphism of the additive groups that also respects multiplication sends invertible elements to invertible elements.
\end{hint}
\begin{hint}{ex:commutator}
$g[a,b]g^{-1}=[gag^{-1},gbg^{-1}]$. In $G/N$, $\pi(a)\pi(b)=\pi(b)\pi(a)$ iff $[a,b]\in N$. For $\mathfrak S_3$, commutators are even.
\end{hint}
\begin{hint}{ex:quaternion}
List the subgroups: each contains $-1$ unless it is trivial. For the last part, count elements of order $2$.
\end{hint}
\begin{hint}{ex:affine}
Composition gives $(a,b)(a',b')=(aa',ab'+b)$. The map $x\mapsto ax+b\mapsto a$ is a homomorphism. For $p=3$, look at the action on the three points.
\end{hint}
\begin{hint}{ex:normalnormal}
For transitivity, look at $\langle s\rangle\le\langle s,sr^2\rangle\le D_4$. For the second part, conjugation by $g$ restricts to an automorphism of $N$.
\end{hint}
\begin{hint}{ex:order8}
Sort by the largest order of an element. If it is $4$, with $a$ of order $4$ and $b\notin\langle a\rangle$, then $bab^{-1}\in\{a,a^{-1}\}$ and $b^2\in\langle a\rangle$.
\end{hint}

\section{Solutions}

\begin{solution}{ex:autzn}
\idea The entire group $\Z/n\Z$ is generated by the residue $[1]$. Start by choosing its image; then check when that choice gives an automorphism rather than merely an endomorphism.

\writeup Let $f:(\Z/n\Z,+)\to(\Z/n\Z,+)$ be a homomorphism and set $a=f([1])$. Every residue is $[k]=k[1]$, for an integer $k$. Repeated addition and the identity $f(-x)=-f(x)$ give
\[
f([k])=k f([1])=[k]a.
\]
This determines $f$ completely. Conversely, for \emph{any} residue $a$, the rule $f_a([k])=a[k]$ is well defined modulo $n$ and satisfies
\[
f_a(x+y)=a(x+y)=ax+ay=f_a(x)+f_a(y).
\]
Thus there is exactly one endomorphism for each residue $a$.

When is $f_a$ invertible? If $a$ has a multiplicative inverse $b$ modulo $n$, then $f_b\circ f_a=f_{ba}=f_1$, so $f_a$ is an automorphism. Conversely, if $f_a$ is onto, there is some residue $b$ with $f_a(b)=1$. This says $ab=1$ modulo $n$, so $a$ is a unit. For this finite group, injectivity and surjectivity are equivalent, but we did not need to assume either without proof.

Finally,
\[
(f_a\circ f_b)(x)=a(bx)=(ab)x=f_{ab}(x).
\]
The assignment $a\mapsto f_a$ therefore respects multiplication, is injective since $f_a(1)=a$, and is surjective onto all automorphisms by the first paragraph. Hence
\[
\operatorname{Aut}(\Z/n\Z)\cong(\Z/n\Z)^*.
\]
For instance, modulo $8$, the four units $1,3,5,7$ give the four automorphisms; multiplication by $2$ is a homomorphism but not an automorphism, since its image has only four residues.
\end{solution}

\begin{solution}{ex:centre}
\idea To prove normality, show that conjugation leaves central elements fixed. For the second statement, lift a generator of the cyclic quotient to $G$ and write every group element as a power of that lift times something central.

\writeup First recall
\[
Z(G)=\{z\in G:zg=gz\text{ for every }g\in G\}.
\]
The identity is in $Z(G)$. If $z,w$ commute with every element, then so does $zw$, because
\[
(zw)g=z(wg)=z(gw)=(zg)w=(gz)w=g(zw).
\]
Also, $zg=gz$ implies $z^{-1}g=gz^{-1}$, so inverses stay in the centre. Hence $Z(G)$ is a subgroup.

For any $g\in G$ and $z\in Z(G)$, centrality gives $gzg^{-1}=z$. Thus every conjugate of every central element is again central. This is precisely $Z(G)\trianglelefteq G$.

Now suppose the quotient $G/Z(G)$ is cyclic, generated by a coset $aZ(G)$. Every $x\in G$ has a coset of the form
\[
xZ(G)=(aZ(G))^m=a^mZ(G)
\]
for some integer $m$. Therefore $x=a^mz$ for some $z\in Z(G)$. Similarly, $y=a^nw$ for some $n\in\Z$ and $w\in Z(G)$.

Powers of the same element $a$ commute. Central elements commute with everything, including one another. Consequently
\[
xy=a^mz\,a^nw=a^{m+n}zw
  =a^{n+m}wz=a^nw\,a^mz=yx.
\]
Since $x$ and $y$ were arbitrary, $G$ is abelian. In that case $Z(G)=G$ and the quotient is actually the trivial cyclic group. In particular, a non-abelian group can never have a non-trivial cyclic quotient by its centre.
\end{solution}

\begin{solution}{ex:phimult}
\idea The Chinese remainder theorem identifies residues modulo $mn$ with pairs of residues modulo $m$ and $n$. Because it preserves multiplication, it also identifies the invertible residues. Count them to recover the multiplicativity of Euler's totient function.

\writeup Assume $\gcd(m,n)=1$. The Chinese remainder theorem gives a bijection
\[
\psi:\Z/mn\Z\longrightarrow\Z/m\Z\times\Z/n\Z,\quad
[k]_{mn}\longmapsto([k]_m,[k]_n).
\]
This is more than a bijection of sets: reduction modulo an integer preserves addition and multiplication. In particular,
\[
\psi([ab]_{mn})=([ab]_m,[ab]_n)
=\psi([a]_{mn})\psi([b]_{mn})
\]
with componentwise multiplication on the right, and $\psi([1])=([1],[1])$.

An element is a unit in a ring if and only if it has a multiplicative inverse. A ring isomorphism sends inverses to inverses, so the units of the product ring correspond exactly to pairs whose two components are units. Explicitly, $[a]_{mn}$ is invertible precisely when both $[a]_m$ and $[a]_n$ are invertible. Restricting $\psi$ to these sets gives a multiplicative group isomorphism
\[
\boxed{(\Z/mn\Z)^*
\cong(\Z/m\Z)^*\times(\Z/n\Z)^*}.
\]

The left-hand group has $\varphi(mn)$ elements, by the definition of Euler's totient. The two groups on the right have $\varphi(m)$ and $\varphi(n)$ elements; their direct product has the product of those numbers. Taking cardinalities gives
\[
\boxed{\varphi(mn)=\varphi(m)\varphi(n)}
\quad\text{whenever }\gcd(m,n)=1.
\]
The coprimality hypothesis is essential: for example, $\varphi(4)=2\ne\varphi(2)\varphi(2)=1$.
\end{solution}

\begin{solution}{ex:commutator}
\idea A conjugate of a commutator is again a commutator, which gives normality. In a quotient, ``two classes commute'' means ``their commutator dies'', which gives the rest.
\writeup For any $g$, $c_g$ is an automorphism, so $g[a,b]g^{-1}=[gag^{-1},gbg^{-1}]$; it maps the generating set of $G'$ to itself, hence $gG'g^{-1}\subseteq G'$. Let $N\trianglelefteq G$ and $\pi:G\to G/N$. Then $\pi(a)\pi(b)=\pi(b)\pi(a)$ iff $\pi([a,b])=e$ iff $[a,b]\in N$. So $G/N$ is abelian iff every commutator lies in $N$ iff $G'\subseteq N$ (as $N$ is a subgroup). With $N=G'$, $G/G'$ is abelian.
In $\mathfrak S_3$, commutators have sign $\varepsilon(a)\varepsilon(b)\varepsilon(a)^{-1}\varepsilon(b)^{-1}=1$, so $\mathfrak S_3'\subseteq A_3$, and $[(1\,2),(1\,3)]=(1\,2)(1\,3)(1\,2)(1\,3)=(1\,2\,3)\ne e$ (compose from the right). As $A_3$ has prime order $3$, $\mathfrak S_3'=A_3$.
\end{solution}

\begin{solution}{ex:quaternion}
\idea Subgroups of a group of order $8$ have order $1,2,4,8$, and $Q_8$ has very few elements of order $2$; that rigidity forces every subgroup to be normal and blocks any decomposition.
\writeup The elements $\pm i,\pm j,\pm k$ have order $4$ (their square is $-1$), and $-1$ has order $2$. So the subgroups are: $\{1\}$; $\{\pm1\}$ (the only one of order $2$); $\langle i\rangle=\{\pm1,\pm i\}$, $\langle j\rangle$, $\langle k\rangle$ of order $4$ (a subgroup of order $4$ contains an element of order $4$, or else all its non-identity elements would have order $2$, but there is only one such element); and $Q_8$. Subgroups of index $2$ are normal, $\{\pm1\}$ is central, and the trivial subgroup and $Q_8$ are normal. But $ij=k\ne-k=ji$.
If $Q_8=N\rtimes H$ with $N,H$ proper and $N\cap H=\{1\}$, both are non-trivial, so both contain $-1$ (the unique element of order $2$, which every non-trivial subgroup contains): a contradiction.
\end{solution}

\begin{solution}{ex:affine}
\idea An affine map is determined by its slope $a$ and its translation term $b$. Compute composition first; the resulting multiplication rule will display both the normal translation subgroup and the complementary scaling subgroup.

\writeup Work in the field $K=\Z/p\Z$. For $a\in K^*$ and $b\in K$, write
\[
f_{a,b}(x)=ax+b.
\]
This map is a bijection, with inverse $f_{a^{-1},-a^{-1}b}$. Two such maps are equal only if they have the same $b=f(0)$ and the same $a=f(1)-f(0)$. There are $p-1$ choices for $a$ and $p$ for $b$, so
\[
|\mathrm{Aff}_p|=p(p-1).
\]
Composition gives
\[
f_{a,b}\circ f_{c,d}(x)
=a(cx+d)+b=acx+(ad+b).
\]
Identifying a map with the pair $(a,b)$, the group law is therefore
\[
(a,b)(c,d)=(ac,ad+b).
\]

The map $\pi:\mathrm{Aff}_p\to K^*$ defined by $\pi(a,b)=a$ is a surjective homomorphism. Its kernel is
\[
T=\{(1,b):b\in K\},
\]
the translations $x\mapsto x+b$. Thus $T\trianglelefteq\mathrm{Aff}_p$ and $T\cong(K,+)\cong\Z/p\Z$.

The maps $H=\{(a,0):a\in K^*\}$ form a subgroup isomorphic to $(K^*,\cdot)$. We have $T\cap H=\{(1,0)\}$, and every affine map factors uniquely as
\[
(a,b)=(1,b)(a,0).
\]
Hence $\mathrm{Aff}_p=TH$ is an internal semidirect product. Its action is visible in conjugation:
\[
(a,0)(1,b)(a^{-1},0)=(1,ab).
\]
So multiplication by $a$ on $\Z/p\Z$ is the relevant action, and
\[
\boxed{\mathrm{Aff}_p\cong\Z/p\Z\rtimes(\Z/p\Z)^*}.
\]

For $p=3$, every affine map is a permutation of the three-element set $K$. The natural map into $\mathfrak S(K)\cong\mathfrak S_3$ is injective because two distinct functions are distinct permutations. Both groups have $3\cdot2=6$ elements, so this injection is onto. Therefore
\[
\mathrm{Aff}_3\cong\mathfrak S_3.
\]
This also gives a geometric way to see the normal three-element subgroup of $\mathfrak S_3$: it is the group of translations of the three residues.
\end{solution}

\begin{solution}{ex:normalnormal}
\idea The lattice of $D_4$ shows a chain $\langle s\rangle\trianglelefteq\langle s,sr^2\rangle\trianglelefteq D_4$ where the bottom subgroup is not normal in $D_4$. For the positive statement, conjugation by $g\in G$ is an automorphism of $N$, and characteristic subgroups survive every automorphism.
\writeup The Klein group $V=\langle s,sr^2\rangle$ has index $2$ in $D_4$, so is normal; $\langle s\rangle$ has index $2$ in $V$, so is normal in $V$; but $r\,s\,r^{-1}=sr^{-2}=sr^2\notin\langle s\rangle$. So normality is not transitive.
If $N\trianglelefteq G$ and $g\in G$, then $c_g$ maps $N$ to $N$ and restricts to an automorphism of $N$. If $K$ is characteristic in $N$, then $c_g(K)=K$, i.e.\ $gKg^{-1}=K$: $K\trianglelefteq G$. In a cyclic group $N$ there is at most one subgroup of each order (Theorem~\ref{thm:cyclic}), and an automorphism maps a subgroup to one of the same order; so every subgroup of $N$ is characteristic, hence normal in $G$.
\end{solution}

\begin{solution}{ex:order8}
\idea Classify by the largest order of an element, which divides $8$. Order $8$ gives the cyclic group; exponent $2$ gives a vector space over $\Z/2\Z$. In the remaining case there is $a$ of order $4$; $\langle a\rangle$ has index $2$, so it is normal, and conjugation by any $b$ outside acts on it as an automorphism of $\Z/4\Z$, which is $a\mapsto a$ or $a\mapsto a^{-1}$ (Exercise~\ref{ex:autzn}). Then $b^2\in\langle a\rangle$, and only two values are possible. Finally the five groups differ in their numbers of elements of each order.
\writeup Let $|G|=8$ and $m$ the largest order of an element, so $m\in\{1,2,4,8\}$ by Lagrange (and $m\ge2$).
\emph{$m=8$}: $G$ is cyclic, $G\cong\Z/8\Z$.
\emph{$m=2$}: $G$ is abelian (Exercise~\ref{ex:square}) and $x+x=0$ in additive notation, so $G$ is a vector space over $\Z/2\Z$ with $8$ elements, of dimension $3$: $G\cong(\Z/2\Z)^3$.
\emph{$m=4$}: let $a$ have order $4$, $N=\langle a\rangle$, and $b\notin N$. $N$ has index $2$, so it is normal, $G=N\cup Nb$, and $b^2\in N$ (in $G/N\cong\Z/2\Z$, $(bN)^2=N$). Conjugation by $b$ is an automorphism of $N\cong\Z/4\Z$, so $bab^{-1}=a$ or $a^{-1}$. Also $b^2$ commutes with $b$, and $b^2\ne a^{\pm1}$ (otherwise $b$ would have order $8$).
\begin{itemize}
\item If $bab^{-1}=a$, $G$ is abelian. If some $b\notin N$ has order $2$, then $N\cap\langle b\rangle=\{e\}$ and $G\cong\Z/4\Z\times\Z/2\Z$ by Proposition~\ref{prop:internal}. Otherwise $b^2=a^2$ for every $b\notin N$, and then $ba^{-1}$ (also outside $N$) has $(ba^{-1})^2=b^2a^{-2}=e$, a contradiction; so this case does not occur.
\item If $bab^{-1}=a^{-1}$, then $G$ is not abelian. Since $b^2\in\{e,a^2\}$: if $b^2=e$, the relations $a^4=b^2=e$, $bab^{-1}=a^{-1}$ are those of $D_4$ (with $r=a$, $s=b$), and the map $r^ks^\varepsilon\mapsto a^kb^\varepsilon$ is an isomorphism (both groups have their $8$ elements written uniquely in this form and multiply by the same rule). If $b^2=a^2$, then $a\mapsto i$, $b\mapsto j$ gives in the same way $G\cong Q_8$ (with $i^4=1$, $j^2=i^2=-1$, $jij^{-1}=-i=i^{-1}$).
\end{itemize}
The five groups are pairwise non-isomorphic: $\Z/8\Z$ has an element of order $8$; $(\Z/2\Z)^3$ has exponent $2$; $\Z/4\Z\times\Z/2\Z$ is abelian of exponent $4$; $D_4$ has $5$ elements of order $2$ and $Q_8$ only one.
\end{solution}

\section*{Sources and further reading}

\begin{itemize}
\item D.~S.~Dummit and R.~M.~Foote, \emph{Abstract Algebra}, 3rd ed., Wiley, 2004, Chapters~3 and~5. Quotient groups, the isomorphism theorems, and semidirect products with many worked examples.
\item M.~Artin, \emph{Algebra}, 2nd ed., Pearson, 2011, Chapter~2. Homomorphisms and quotients, with the correspondence theorem.
\item J.~J.~Rotman, \emph{An Introduction to the Theory of Groups}, 4th ed., Springer, 1995, Chapters~2 and~7. A deeper treatment, including extensions and the classification of small groups.
\item K.~Conrad, \emph{Groups of order $8$}, and other expository notes on semidirect products, on the author's web page at the University of Connecticut.
\item D.~Joyner, \emph{Adventures in Group Theory}, 2nd ed., Johns Hopkins University Press, 2008. Parity and the legal positions of Rubik's cube.
\end{itemize}
